\section{From finite rectangles to certified geometric bounds}\label{sec:certificates}

The analytic reduction needs uniform inequalities on entire root
regions and on all line directions. A numerical value attached to a
sampled root is insufficient. This section describes a finite
rational certificate contract and proves its mathematical soundness.
Whether every concrete certificate has been supplied and checked is
a separate question, addressed in Section~\ref{sec:status}.

\subsection{A uniform chord envelope on a closed root box}
Let \(F\) be a nonnegative rational rectangular step field bounded by
\(M\), with zero outside the unit square. For direction \((1,u)\),
\(u>0\), write
\[
 C(v;u)=\int_\R F(t,v+ut)\dd t.
\]
For a closed root box \(Q_0=[a_0,a_1]\times[b_0,b_1]\), every root
has intercept in
\(l(u)=b_0-u a_1\le v\le h(u)=b_1-u a_0\).
For now assume \(Q_0\) is not a single point.

At each grid vertex \((a,\eta)\), with exterior cells read as zero,
define the signed corner coefficient
\[
 k_{a,\eta}=F_{\rm LU}+F_{\rm RD}-F_{\rm RU}-F_{\rm LD}.
\]
\begin{lemma}[Corner representation]\label{lem:corners}
For \(u>0\),
\begin{equation}\label{eq:corners}
 C(v;u)=\frac1u\sum_{a,\eta} k_{a,\eta}(v-\eta+ua)_+.
\end{equation}
\end{lemma}
\begin{proof}
For the indicator of \((a,b)\times(c,d)\), integrate
\(\ind_{v+ut>c}-\ind_{v+ut>d}\) over \(a<t<b\).
The first integral is
\([(v+ub-c)_+-(v+ua-c)_+]/u\), with the analogous expression
for \(d\). The four corner coefficients are \(-1,+1,+1,-1\).
Add these identities with the rational cell weights and combine
shared vertices. Nonaxis chords meet the zeroed grid boundaries
in finitely many points, so those values do not change the integral.
\end{proof}

For a vertex with negative coefficient set
\[
 P=a_1-a+\frac{\eta-b_0}{u},\qquad
 Q=a-a_0+\frac{b_1-\eta}{u},\qquad
 T_{a,\eta}(u)=\max\left(0,\min\left(P,Q,\frac{P+Q}{4}\right)\right).
\]
Here \(Q\) is a scalar; the root box is \(Q_0\).
Define
\begin{align}
 U_1(u)&=\max(C(l;u),C(h;u))+
             \sum_{k_{a,\eta}<0}(-k_{a,\eta})T_{a,\eta}(u),
             \label{eq:U1}\\
 U_2(u)&=C(l;u)+\sum_{J>0}J\,
       \big|[a,b]\cap[(\eta-h)/u,(\eta-l)/u]\big|.
       \label{eq:U2}
\end{align}
In \eqref{eq:U2} the sum runs over horizontal grid edges
\((a,b)\times\{\eta\}\) with upward jump
\(J=F_{\rm up}-F_{\rm down}>0\).

\begin{proposition}[Closed-box envelope]\label{prop:envelope}
For every root \((x,y)\in Q_0\), \(v=y-ux\), and \(u>0\),
\begin{equation}\label{eq:envelope}
 0\le C(v;u)\le U(u):=\min(U_1(u),U_2(u),M)\le M.
\end{equation}
The envelope is a finite expression formed from rational functions
\(\alpha+\beta/u\), rational arithmetic, minimum, maximum and
positive part.
\end{proposition}
\begin{proof}
For a corner term \(k(v-v_0)_+/u\), compare it with its endpoint
linear interpolation on \([l,h]\). A nonnegative \(k\) has no
positive excess. A negative \(k\) has excess zero if the corner
lies outside the interval, and otherwise maximum
\[
 (-k)\frac{(v_0-l)(h-v_0)}{u(h-l)}
 =(-k)\frac{PQ}{P+Q}
 \le(-k)\min(P,Q,(P+Q)/4).
\]
Sum the corner identities. Interpolation commutes with finite
addition, and the interpolated total is at most its larger endpoint.
This gives \(U_1\).

For fixed horizontal coordinate \(t\), move vertically from
\(l+ut\) to \(v+ut\). The field difference is the sum of its
crossed horizontal jumps. Discard negative jumps and extend the
upper endpoint to \(h+ut\). Integrating the resulting indicator
inequality over \(t\) gives exactly \eqref{eq:U2}. Boundary
exceptions involve finitely many values of \(t\). The trivial
bound \(C\le M\) follows because the major-coordinate length
inside the square is at most one. Each of the three bounds is
nonnegative. The displayed expressions for the intersections and
corners have the asserted finite form.
\end{proof}

For a point root, simply use \(C(l;u)\). A root box of zero width
or zero height is covered by the same argument whenever \(h-l>0\).
At \(u=0\) one may set \(U(0)=M\) for the angular integral.
This does not remove the separate obligation to check axis capacities.

\subsection{Verifying branches without trusting the branch generator}
Suppose a proposed branch on a rational interval \([r,s]\subset(0,1]\)
is \(V(u)=a+b/u\). To compare it with
\(a'+b'/u\), multiply the difference by positive \(u\).
Its sign is the sign of the affine function
\((a-a')u+(b-b')\), and is therefore determined by its two endpoint
values. Induction on the finite expression tree proves every min/max
branch choice throughout the interval. Check that the proposed
intervals cover \([0,1]\) in order, with no gaps in their interiors.

A finite exact branch partition exists: on each common child
interval, two nonidentical branches cross at at most one rational
point. Add that point if it lies inside the interval; identical
branches require no subdivision. Near zero, the final bounded
nonnegative envelope must have \(b=0\) on its first branch. A checker
can either verify this or use the constant bound \(M\) on that first
interval. It need not trust the search procedure that supplied the
partition.

\subsection{Rational integral bounds without logarithms}
On \(0<r<s\), suppose \(V(u)=a+b/u\) has nonnegative endpoints.
It is monotone, hence nonnegative everywhere on the interval.
Let
\[
 H_E=(s-r)\max(V(r)^2,V(s)^2).
\]
When \(b\ge0\), take the secant line
\[
 L(u)=a+b\frac{r+s-u}{rs}.
\]
When \(b<0\), take \(c=(r+s)/2\) and
\[
 L(u)=a+2b/c-bu/c^2.
\]
\begin{lemma}[Affine integral certificate]\label{lem:affine}
In both cases \(0\le V\le L\), and with
\(\alpha=L(r)\), \(\beta=L(s)\),
\begin{equation}\label{eq:affine}
 \int_r^sV(u)^2\dd u
 \le\min\left(H_E,\frac{s-r}{3}
                       (\alpha^2+\alpha\beta+\beta^2)\right).
\end{equation}
Every quantity on the right is rational.
\end{lemma}
\begin{proof}
For \(b\ge0\),
\[
 \frac{r+s-u}{rs}-\frac1u=\frac{(u-r)(s-u)}{rsu}\ge0.
\]
For \(b<0\),
\[
 L(u)-V(u)=\frac{(-b)(u-c)^2}{uc^2}\ge0.
\]
Square the inequalities, and integrate the square of the affine
function using \(u=r+t(s-r)\). Its integral is
\((s-r)(\alpha^2+\alpha\beta+\beta^2)/3\).
The endpoint bound \(H_E\) is also valid by monotonicity. The smaller
of two upper bounds for the same integral remains an upper bound.
\end{proof}

If the first interval is \([0,s]\) with verified \(b=0\), its
integral is \(a^2s\); no division by zero is performed.
At fixed integer scale \(S>0\), replace an interval upper bound
\(H\) by \(\lceil SH\rceil/S\). The error is less than \(1/S\)
per interval. At a fixed partition and scale,
\(\lceil S\min(H_E,H_A)\rceil\le\lceil SH_E\rceil\), so this
improvement cannot worsen an existing endpoint certificate.

A strict integral margin guarantees some finite partition and
sufficiently fine rounding scale will succeed. It does not guarantee
that indefinitely subdividing at a fixed scale succeeds: the
accumulated rounding budget depends on the number of intervals.
Nor does a handful of successful leaves certify the remaining leaves.

Apply the envelope and integral bounds to the four direction families
and each of the three other types. Reflections and coordinate exchange
preserve the major-coordinate parameter, so no angular or arc-length
Jacobian is inserted. If the resulting twelve integral upper bounds
sum to \(H^+\), then throughout the root box
\begin{equation}\label{eq:certRisk}
 \mathcal R_r(F;x)\le\frac{D_*}{24}H^+.
\end{equation}
For unscaled fields followed by time contraction by \(s_0\), the
coefficient is \(D_*s_0^2/24\). Reusing a bound on symmetry images
requires verified invariance of the relevant deterministic fields,
or separately transformed fields. Symmetry of the random output is
neither assumed nor required.

\subsection{Capacity: axes, nonaxes and the infinite tail}
For a fixed nonaxis slope \(q\ne0\), the chord of a cell
\((A,B)\times(C,D)\) at intercept \(v\) is
\[
 \pos{\min\left(B,\max\left(\frac{C-v}{q},\frac{D-v}{q}\right)\right)
 -\max\left(A,\min\left(\frac{C-v}{q},\frac{D-v}{q}\right)\right)}.
\]
It is continuous and piecewise affine in \(v\), with switches among
\(C-qA,C-qB,D-qA,D-qB\). Summing finitely many cell contributions
proves that the maximum capacity is attained at an intercept through
a grid vertex. Thus a finite rational list covers every intercept.

For an axis the assertion about endpoint values is false. A horizontal
line in an open row strip has chord equal to that strip's weighted
row sum, whereas a horizontal grid boundary may have chord zero.
Check all \emph{open strip values} and the actual boundary value zero;
use column sums for the vertical axis. Finer zeroed grid lines either
remove finitely many points of a nonaxis chord or set an entire axis
chord to zero, so a coarse nonnegative envelope remains an upper bound.

For the candidate's height cutoff, all directions with \(h\ge8\)
satisfy
\[
 L_{v,r}(F;x)\le\frac{\|F_r+F_{r+\bar v}\|_\infty}{2h}
 \le\frac{\|F_r+F_{r+\bar v}\|_\infty}{16}.
\]
Check the finite cell maxima for all six unordered type pairs, and
check each primitive direction of heights one through seven by the
appropriate axis or nonaxis rule. Nonnegative accumulation makes a
terminal-field capacity bound valid at earlier times as well.

\begin{proposition}[A finite geometric checking contract]\label{prop:contract}
A successful finite checker must establish all of the following:
\begin{enumerate}
\item The exact field values, boundary convention, phase durations,
types, order, and accepted union are the specified inputs.
The actual accumulated field is bounded by the corresponding
reference prefix at every used phase.
\item Every used root lies in a checked closed box; all required
direction families and type pairs have complete angular coverage.
Every branch is justified from its field and box, and every integral
bound uses sound rational arithmetic and outward rounding.
\item The capacities cover the axis strips, the finite nonaxis list,
and the infinite-height tail. The union area counts overlaps once.
\item The returned rational bounds satisfy
\(U_H<K\), \(U_L<\lambda<1-\rho\), and \(U_R<d<1\).
\end{enumerate}
Then the analytic input inequalities hold with uniform positive
margins \(K-U_H\), \(\lambda-U_L\), and \(d-U_R\).
\end{proposition}
\begin{proof}
The capacity rules cover all intercepts and directions.
Proposition~\ref{prop:envelope}, branch induction and
Lemma~\ref{lem:affine} imply \eqref{eq:certRisk} on each closed box.
Nonnegative field domination transfers this to the actual field;
the phase and root-deadline matching ensures it is the correct
field. The finite rational maxima or a common verified threshold
give one uniform bound, whose positive difference from the target
is independent of the grid size and root. The separate union-area
calculation supplies the mass. Proposition~\ref{prop:pair} then
transfers the continuum risk to the discrete pair sum.
\end{proof}

Notice that a stored numerical upper bound or a hash of a data file
does not by itself establish any of the branch, coverage or prefix
premises. Hashes identify the input being discussed; the mathematical
checking contract establishes the inequalities about that input.

